% BEGIN JE20261104_EXTENSION
\clearpage
\section{Fee-adjustment counterfactual and regular-distribution extension}
\label{app:extension}

\subsection{Proof of Corollary 1: frozen versus reoptimised fees}
Fix any $\lambda_0$ in Proposition~1's compact feasible interval and
hold the private-only fee at
$f_i^P(\lambda_0)=[U+D(\lambda_0)]/2$, while varying
$\lambda$ in a neighbourhood of $\lambda_0$. Because $U$ is independent
of $\lambda$, the counterfactual participation cutoff is
\[
 \ell^{\mathrm{fr}}(\lambda)
 =\frac{f_i^P(\lambda_0)-U}{d+\eta_i\lambda},\qquad
 \left.\frac{d\ell^{\mathrm{fr}}}{d\lambda}\right|_{\lambda_0}
 =-\frac{\eta_i\ell(\lambda_0)}{D(\lambda_0)}<0.
\]
The alternative public-basic-plus-premium continuation has universal
participation and premium fee $d/2$, independent of $\lambda$.
The fixed-fee public-cost ceiling is obtained from
\eqref{eq:screen-threshold} by replacing the private-only optimum
$\ell(\lambda)$ with $\ell^{\mathrm{fr}}(\lambda)$. At $\lambda_0$,
\begin{equation}
 \left.\frac{dT_i^{\mathrm{fr}}}{d\lambda}\right|_{\lambda_0}
 =-\frac{\eta_i\ell}{D}
  \left[N_i(u_i+d\ell)+C\left(S_i^0+
              \frac{\lambda_0K_i\ell}{2}\right)\right]
    +\frac{CK_i\ell^2}{4},
 \label{eq:fr-derivative-proof}
\end{equation}
where $\ell=\ell(\lambda_0)$ and $D=D(\lambda_0)$.
Under Proposition~1, $K_i<0$, $\eta_i>0$, and $\ell>0$.
The bracket equals $N_i(u_i+d\ell_0)+O(\gamma)>0$ uniformly
on the compact interval for sufficiently small positive $\gamma$.
Both terms of equation~\eqref{eq:fr-derivative-proof} are therefore
strictly negative, whereas Proposition~1 proves the positive
reoptimised derivative. This compares the \emph{same} objective under
two private-fee adjustment rules, not two fee equilibria.
For the original exact Cournot/rational baseline,
$T_B(0)=0.0404250710071$, $T_B(0.5)=0.0405316701024$;
freezing $f_B^P$ at its $\lambda=0$ optimum instead gives
$T_B^{\mathrm{fr}}(0.5)=0.0400020105974$.
All three numbers are reproduced from the original primitives by
\texttt{je\_extension\_cert.py}; none is a new fitted calibration.

\subsection{Proof of Proposition 2: general founder ability distribution}
Let $m(t)=1-F_i(t)$, $q(t)=\int_t^1xg_i(x)\,dx$, and
$J(t)=\int_0^t xg_i(x)\,dx$. The partner's full-coverage moment is
$\bar\theta_j$. Along the private-only continuation write
\[
 U=u_i+\gamma N_j\omega_iR_0,
 \qquad D=d+\gamma N_j\omega_i\Pi\bar\theta_j\lambda.
\]
Any strictly interior posted fee implements a cutoff $\ell$
through $f_i=U+D\ell$. The private intermediary's gross receipts are
$N_i(U+D\ell)m(\ell)$ and the first-order condition is
\begin{equation}
 Dm(\ell)=(U+D\ell)g_i(\ell).
 \label{eq:general-fee-foc}
\end{equation}
The hazard assumption yields a global optimum rather than merely a
stationary point. In fact the derivative of revenue with respect to
$\ell$ has the sign of
$1-[\ell+U/D]h_i(\ell)$; this expression decreases
strictly through zero when $h_i$ is nondecreasing and positive.
The stated inequality $u_i h_i(0+)<d$ ensures a positive interior
cutoff at $\gamma=0$, and the hazard must become sufficiently large
near the upper endpoint of a compact support with full mass.
The unique strict optimum persists uniformly for sufficiently small
$\gamma$ over the feasible compact $\lambda$ interval.
In co-provision, all founders participate and the premium cutoff
$\tau_B$ globally maximises $d\tau[1-F_i(\tau)]$;
it solves $\tau_Bh_i(\tau_B)=1$ and is independent of $\lambda$.

The gross regional public-addition threshold is exactly
\begin{equation}
 \begin{split}
 T_i^F(\lambda)={}&N_i\bigl\{u_iF_i(\ell)
             +d[q(\tau_B)-q(\ell)]\bigr\}\\
 &+C\bigl\{S_i^0F_i(\ell)
          +\lambda K_i\bar\theta_j J(\ell)\bigr\}.
 \end{split}\label{eq:general-ceiling-exact}
\end{equation}
Differentiating the private optimum condition
\eqref{eq:general-fee-foc} at fixed $U$ gives
\begin{equation}
 \frac{\partial\ell}{\partial D}
 =\frac{U g_i(\ell)}{D[2Dg_i(\ell)+(U+D\ell)g_i'(\ell)]}>0,
 \quad
 \ell'=\gamma N_j\omega_i\Pi\bar\theta_j
              \frac{\partial\ell}{\partial D}.
 \label{eq:general-cutoff-response}
\end{equation}
The denominator is strictly positive at the unique
hazard-based global maximum. The welfare derivative, retaining the
endogenous cutoff response rather than applying the seller's envelope
to government welfare, is exactly
\begin{equation}
 \begin{split}
 \frac{dT_i^F}{d\lambda}={}&
 \bigl\{N_i(u_i+d\ell)+
   C[S_i^0+\lambda K_i\bar\theta_j\ell]\bigr\}
     g_i(\ell)\ell'\\
 &+CK_i\bar\theta_j J(\ell).
 \end{split}\label{eq:general-slope-exact}
\end{equation}
At zero matching intensity the unique optimal cutoff is
$\ell_0$. A Taylor expansion about $\gamma=0$ gives
$\ell=\ell_0+O(\gamma)$ and
\[
 \ell'=\gamma N_j\omega_i\Pi\bar\theta_j
   \frac{u_i g_i(\ell_0)}{d B_F}+O(\gamma^2).
\]
Substitute this into equation~\eqref{eq:general-slope-exact}, use
$C=\gamma N_iN_j$, and collect order-$\gamma$ terms to obtain
equation~\eqref{eq:general-slope}. The remainder is uniform because
the implicit optimum has a uniformly positive second-order
denominator on a neighbourhood of $\ell_0$ and $g_i$ is twice
continuously differentiable there. The strict lower band on
$\omega_i$ makes the leading coefficient positive, while the upper
band gives $K_i<0$. As in Proposition~1, $p_b>0$ then implies
$S_i^0=p_bK_i-\omega_i\zeta^2\Pi^2/(2k)<0$.
The zero-matching limit of equation~\eqref{eq:general-ceiling-exact}
is the direct-service term stated in Proposition~2;
its strict positivity, if imposed, persists for sufficiently small
positive $\gamma$. For $F_i(t)=t$,
$\ell_0=(d-u_i)/(2d)$, $J_F=\ell_0^2/2$,
$A_F=u_i(u_i+d)/(4d^2)$, and $\bar\theta_j=1/2$,
which reduce the result identically to equation~\eqref{eq:main-primitive-slope}.

\subsection{Exact-rational finite-intensity Beta$(2,2)$ sign certificate}
This separate test checks the \emph{conditional} region-$B$ threshold
at the original parameter values $\gamma=0.1$, $\omega_B=0.2$,
$\lambda\in[0,1/2]$, while region $A$ remains fully participating.
For the Beta$(2,2)$ law,
$F(t)=3t^2-2t^3$, $g(t)=6t(1-t)$,
$J(t)=2t^3-3t^4/2$, and
\[
 h'(t)=\frac{6(1+2t^2)}{(1-t)^2(1+2t)^2}>0.
\]
The optimal private cutoff is the unique interior root of
$8D\ell^2+(6U-D)\ell-D=0$. The co-provision premium
cutoff is $\tau_B=(1+\sqrt{33})/16$, independent of $\lambda$.
With the frozen rational Cournot and research primitives, define
$D_{\max}=d+\eta_i/2$. Exact fractions verify the first-order-condition
signs at $\ell=3/10$ with $D=d$ and at $\ell=33/100$ with
$D=D_{\max}$, establishing
\[
 \frac3{10}<\ell(\lambda)<\frac{33}{100}
 \quad\text{for every }\lambda\in[0,1/2].
\]
Because $g$ increases and $g'$ decreases on this interval, equation
\eqref{eq:general-cutoff-response} supplies the exact rational lower bound
\[
 \ell'\geq
 \frac{\eta_i U g(3/10)}{
 D_{\max}\{2D_{\max}g(33/100)
       +[U+D_{\max}(33/100)]g'(3/10)\}}
 =:\underline{\ell'}>0.
\]
Since $K_B<0$ and $S_B^0<0$, the first brace in
\eqref{eq:general-slope-exact} is bounded from below by
$B_{\min}=N_B(u_B+3d/10)+C S_B^0+C K_B(33/100)/4>0$.
Finally, exact rational interval arithmetic gives
\[
 \frac{dT_B^{\mathrm{Beta}}}{d\lambda}
 \geq B_{\min}g(3/10)\underline{\ell'}
       +\frac{CK_B}{2}J(33/100)
 >\frac1{10000}>0
\]
throughout $\lambda\in[0,1/2]$.
The last inequality is a conservative certified bound, not a floating
point extrapolation between endpoints. The independently reexecuted
\texttt{je\_extension\_cert.py} contains the exact
\texttt{Fraction} checks. It neither re-proves all coupled
$(P,P)$ fee equilibria nor classifies all sixteen entry
configurations for the Beta alternative. Thus the conditional positive
slope is portable to this particular distribution, while the
reported four-party ownership switch is not.
% END JE20261104_EXTENSION
